\documentclass[twocolumn]{aastex631}
\begin{document}
\title{The reionization ionizing-photon-budget shortfall for star-forming galaxies is not robust to systematics at z\~{}5 (highC systematic corner)}
\author{NebulaMind Lab (autonomous pipeline)}
\affiliation{NebulaMind Lab, \url{https://lab.nebulamind.net}}
\begin{abstract}
Reionization ionizing-photon-budget reconciliation at z\~{}5: star-forming galaxies require f\_esc=0.036 (+0.034/-0.018) to close the budget (Madau-Dickinson SFRD, log xi\_ion=25.5+/-0.15, clumping C=2-5, JWST-SFRD tail), versus indirect-proxy-inferred f\_esc=0.062 (+0.108/-0.039) from LzLCS O32/beta calibrations. Median delta(required-inferred)=-0.023 dex-frac (16-84\%: -0.131 to +0.023); 32\% of the systematic MC shows a shortfall. Conclusion: the budget CLOSES within the systematic, and the sign holds under both O32 and beta calibrations. The result is bounded by the xi\_ion x clumping x proxy-calibration systematic, not by statistics. Generated autonomously from public data (jwst) via the NebulaMind Lab runner: a bounded, reproducible, descriptive study.
\end{abstract}
\section{Introduction}
Recent studies have highlighted potential discrepancies in the ionizing photon budget during the reionization epoch, sparking concerns about our understanding of this critical period in cosmic history [Muñoz2024]. In particular, there is a need to reconcile the estimated ionizing photon production from star-forming galaxies with the required amount to maintain reionization. This issue has been explored through various approaches, including excursion set models [Park2022] and assessments of galaxy ionizing photon budgets at high redshifts [Duncan2015]. However, these efforts have not yet fully resolved the tension between observed and required ionizing photon rates.
\section{Data and method}
To address this challenge, we adopt a literature-anchored budget calculation that utilizes established values from previous research. Specifically, we employ the cosmic star formation rate density (SFRD) analytic fitting function proposed by Madau \& Dickinson (2014), along with published calibrations for the ionizing photon production efficiency (xi\_ion) and escape fraction (f\_esc) proxies [LzLCS: Chisholm+22, Flury+22; Simmonds+24]. Our method focuses on reconciling these values to determine if star-forming galaxies can account for the necessary ionizing photons during reionization.
\section{Result}
Our analysis reveals that at z\~{}5, star-forming galaxies require an escape fraction of f\_esc=0.036 (+0.034/-0.018) to close the ionizing photon budget. This value is compared to indirect proxy-inferred estimates of f\_esc=0.062 (+0.108/-0.039) derived from LzLCS O32/beta calibrations. The median difference between these values is -0.023 dex-frac, with 16-84\% confidence intervals ranging from -0.131 to +0.023. Notably, 32\% of our systematic Monte Carlo simulations indicate a shortfall in the ionizing photon budget.
\begin{figure}\centering\includegraphics[width=\columnwidth]{result.png}\caption{The reionization ionizing-photon-budget shortfall for star-forming galaxies is not robust to systematics at z\~{}5 (highC systematic corner)}\end{figure}
\section{Caveats}
It is essential to acknowledge that our approach relies on automated, single-selection, and uncalibrated measurements, which introduces limitations to our findings. The accuracy of our results depends heavily on the assumptions underlying the adopted calibrations and the Madau \& Dickinson (2014) SFRD fitting function. Furthermore, our analysis does not account for potential systematic uncertainties in the observational data or the impact of other factors that may influence the ionizing photon budget, such as variations in galaxy properties or environmental effects. These caveats highlight the need for continued research and refinement of our understanding of reionization dynamics.
\section*{References}
\footnotesize
[Muoz2024] Muñoz et al. (2024). Reionization after JWST: a photon budget crisis?. bibcode 2024MNRAS.535L..37M \par [Park2022] Park et al. (2022). Calibrating excursion set reionization models to approximately conserve ionizing photons. bibcode 2022MNRAS.517..192P \par [Duncan2015] Duncan et al. (2015). Powering reionization: assessing the galaxy ionizing photon budget at z < 10. bibcode 2015MNRAS.451.2030D \par [Davies2021] Davies et al. (2021). The Predicament of Absorption-dominated Reionization: Increased Demands on Ionizing Sources. bibcode 2021ApJ...918L..35D \par [Madau2017] Madau (2017). Cosmic Reionization after Planck and before JWST: An Analytic Approach. bibcode 2017ApJ...851...50M

\end{document}
